\section{What Drives GATE-KD?}
\label{sec:analysis}

We organize the analysis around two empirical observations. First, within the
evaluated controls, the endpoint interface obtains most of its numerical gain
after establishing fitted coordinate comparability;
the particular fixed subspace and repeated gauge updates are secondary design
choices. Second, native-space $H_0$ provides a smaller complementary signal
rather than replacing sample-wise endpoint supervision. Unless stated
otherwise, controls use configuration (c), whose 384-dimensional student has
the strongest width bottleneck. The variants come from matched rerun suites;
their GATE-KD reference rows repeat configuration (c) from the main sweep for
consistency. Extended controls are grouped by claim in
Appendices~\ref{app:endpoint-analysis} and~\ref{app:structural-analysis}.

\subsection{Fitted Coordinate Comparability Explains Most of the Endpoint Gain}
\label{sec:analysis-interface}

Table~\ref{tab:gauge-mechanism} holds the PCA subspace and full GATE-KD
objective (endpoint supervision plus native-space $H_0$) fixed while varying
how the endpoint coordinates are treated. The unaligned row trains directly
against the fixed PCA coordinates and therefore doubles as the matched
PCA-only endpoint baseline, the sentence-embedding analogue of the TCS target
construction \citep{tcs}. In configuration (c), no alignment,
a random orientation, a fixed fitted gauge, one refit, and epoch-wise
Procrustes reach 72.51, 73.02, 74.63, 74.64, and 74.79 Avg., respectively. The
additional invariant control directly minimizes the endpoint loss over an
orthogonal rotation on every minibatch and reaches $72.74\pm0.07$. It improves
only marginally over the unaligned controls and trails epoch-wise Procrustes by
2.05 points. Separately, within the gauge-schedule controls, the fixed fitted
gauge gains 2.12 points over no alignment and retains all but 0.16 points of the
epoch-wise result. Taken together, these results argue against explaining the
gain by merely weakening pointwise matching through a student-adaptive
rotation. They are instead consistent with the benefit coming from one fitted
orientation, shared across minibatches, that is consistent with the student's
initial coordinates. We note what these controls do not isolate: because
$R^{(0)}$ is fitted to the initialized student, the schedule controls cannot
separate a resolved cross-model coordinate mismatch from a target orientation
that merely makes early optimization easier for that particular
initialization. A control that fits $R^{(0)}$ on one initialization and trains
independently initialized students against that fixed orientation would
separate the two readings; we leave it to future work and interpret the gauge
here as an empirically useful optimization device rather than as a
transferable coordinate system.

% Control rows: review_h0_pca_controls_20260906-143901, gauge family.
% Epoch-wise Procrustes reference: main-results configuration (c).
% Invariant row: invariant_procrustes_c_3seeds_20260906-215912.
\begin{table}[!t]
    \centering
    \caption{Compact orientation treatments under the full objective (mean
    $\pm$ sample standard deviation, three seeds).}
    \label{tab:gauge-mechanism}
    \small
    \newcommand{\gcell}[2]{#1{\tiny $\pm$ #2}}
    \setlength{\tabcolsep}{5pt}
    \begin{tabular}{@{}lcc@{}}
        \toprule
        Orientation treatment & Updates & Config. (c) Avg. $\uparrow$ \\
        \midrule
        None (fixed PCA target) & 0 & \gcell{72.51}{0.05} \\
        Random orthogonal & 0 & \gcell{73.02}{0.05} \\
        Per-minibatch invariant & Every step & \gcell{72.74}{0.07} \\
        Fixed $R^{(0)}$ & 0 & \gcell{74.63}{0.07} \\
        One refit & 1 & \gcell{74.64}{0.06} \\
        Epoch-wise Procrustes & 4
        & \gcell{\textbf{74.79}}{0.06} \\
        \bottomrule
    \end{tabular}
\end{table}


Within these controls, most of the observed benefit therefore accompanies an
initially fitted gauge: the one-refit and fixed-gauge means differ by only 0.01
points, while epoch-wise refitting adds a modest 0.16 points over the
fixed-gauge result. We interpret refitting as a secondary optimization
refinement, not as an independently established generalization mechanism. The
invariant control uses batches of 128 in a 384-dimensional student space, so
its rotation is underdetermined outside the current batch span. It tests the
reviewed local invariant relaxation but does not rule out every possible global
rotation-invariant objective. Matched subspace controls further show that
aligned random projections remain within 0.13--0.34 points of PCA despite
retaining substantially less teacher energy, while the updated uncentered-SVD
run is 0.19 points below PCA. Within these controls, the proximity of the
subspace constructions is more consistent with fitted coordinate comparability
than with retained variance or a uniquely optimal PCA basis being the primary
driver. Comparisons with learned endpoint maps, subspace-construction controls,
target-rank sensitivity, and gauge-fit-set size appear in
Appendix~\ref{app:endpoint-analysis}.

\FloatBarrier

\subsection{\texorpdfstring{$H_0$}{H0} Provides a Complementary
Connectivity-Scale Signal}
\label{sec:analysis-structural}

In the structural-control comparison, endpoint-only supervision reaches
74.26 Avg., while native-space $H_0$ at
$\lambda_{\mathrm{topo}}=0.5$ reaches 74.79, a 0.53-point refinement.
Table~\ref{tab:structural-controls} compares it with pairwise, local-scale,
teacher-MST, and Gram objectives while holding the endpoint interface and
optimization protocol fixed. The ``None'' row repeats the endpoint-only
reference from Table~\ref{tab:signal-decomposition}, while the $H_0$ row repeats
configuration (c) from the main comparison.

``None'' denotes endpoint-only supervision. The alternatives test a sorted
spanning path, directed nearest-neighbor distances, teacher-MST edge
correspondence, and dense native-space Gram matching. The $H_0$ objective
instead matches the sorted MST edge-weight distributions without edge
correspondence. Appendix~\ref{app:structural-control-definitions} gives the
complete definitions and constraint counts.

% Structural alternatives: latest structural-control runs supplied by the author.
% None reference: Table 8 endpoint-only row. H0 reference: main-results config. (c).
\begin{table}[H]
    \centering
    \caption{Matched structural controls in configuration (c).
    Mean $\pm$ sample standard deviation over three seeds; best value is bold.}
    \label{tab:structural-controls}
    \small
    \newcommand{\scell}[2]{#1{\tiny $\pm$ #2}}
    \setlength{\tabcolsep}{4pt}
    \begin{tabular}{@{}lcc@{}}
        \toprule
        Auxiliary term & Avg. $\uparrow$ & $\Delta$ endpoint \\
        \midrule
        None & \scell{74.26}{0.02} & --- \\
        Sorted spanning path & \scell{73.66}{0.04} & $-0.60$ \\
        kNN distribution & \scell{{74.47}}{0.01} & $+0.21$ \\
        Teacher-MST edges & \scell{74.33}{0.08} & $+0.07$ \\
        Native Gram matching & \scell{73.57}{0.01} & $-0.69$ \\
        $H_0$ & \scell{\textbf{74.79}}{0.06} & $+0.53$ \\
        \bottomrule
    \end{tabular}
\end{table}


Among the evaluated controls and weights, the $H_0$ objective is the strongest
auxiliary, exceeding kNN
distribution matching by 0.32 points and teacher-MST matching by 0.46 points;
sorted spanning-path and native Gram matching instead reduce Avg. by 0.60 and
0.69 points. Thus the gain is not explained by adding an arbitrary structural
objective and is consistent with connectivity-scale matching providing a
distinct complementary signal under the evaluated weights. Appendix
\ref{app:signal-decomposition} shows that $H_0$ alone reaches 64.22, so
connectivity does not replace pointwise correspondence. Topology-weight
sensitivity appears in Appendix~\ref{app:topology-sensitivity}.

\FloatBarrier
